\input{preamble}
\begin{document}

\title{Campinas--Miami\\
Mathematical Meeting 2026}
\maketitle
\phantomsection
\label{section-phantom}
\hfill
\href{http://github.com/danimalabares/stack}
{github.com/danimalabares/stack}
\tableofcontents

\medskip\noindent
{\bf Conference information.}

\href{%
https://sites.google.com/unicamp.br/%
campinas-miami-math-2026/%
}{Campinas--Miami Mathematical Meeting 2026}

May 11--13, IMECC/Unicamp,
Campinas-SP.

\medskip\noindent
{\bf Organizing Committee.}

Leonardo Cavenaghi, Lino Grama,
Ludmil Katzarkov, Jaqueline Mesquita,
Ricardo Miranda, Misha Verbitsky.

\section{Collapsing Flat
SU(2)-Bundles to
Spherical 3-Manifolds}
\label{section-spherical-3-manifolds}

\medskip\noindent
{\bf Speaker.}

Eder M. Correa
(Campinas, Brazil)

\medskip\noindent
{\bf Title.}

Collapsing Flat
SU(2)-Bundles to
Spherical 3-Manifolds.

\medskip\noindent
{\bf Abstract.}

In this talk, we present a geometric
mechanism for the emergence of
spherical $3$-manifolds from the
superspace of Riemannian metrics
associated with flat
SU(2)-bundles over closed orientable
hyperbolic surfaces. We show that any
spherical $3$-manifold can be realized
as a boundary point in the
Gromov-Hausdorff closure of this
superspace through a collapsing
process. Furthermore, we discuss how
the convergence toward the spherical
limit is controlled by the order of
the fundamental group of the target
manifold and the systolic geometry of
the hyperbolic base. Finally, we
illustrate these results by showing
how the Bolza surface provides the
optimal error bound for the
convergence toward the Poincar\'e
homology sphere.

\medskip\noindent
{\bf Main theme.}

Spherical $3$-manifolds appear as
Gromov--Hausdorff boundary points of
superspaces of suitable
$\operatorname{SU}(2)$-bundles over
closed hyperbolic surfaces.

\medskip\noindent
{\bf Spherical manifolds.}

\begin{definition}
A Riemannian manifold $M^n$ is
{\it spherical} if it is complete and
locally isometric to the round sphere
$S^n$.
Equivalently, it has constant
sectional curvature $+1$ and
universal cover $S^n$.
\end{definition}

If $M$ is connected, then
$$
        M \cong S^n/\Gamma
$$
where $\Gamma<\operatorname{Isom}(S^n)$
is finite and acts freely.

In dimension $3$ we use
$$
        S^3\cong \operatorname{SU}(2),
$$
via the unit quaternions.
Thus one gets spherical quotients
$$
        \operatorname{SU}(2)/\Gamma
$$
when $\Gamma<\operatorname{SU}(2)$ is
finite and acts by left multiplication.

The finite subgroups of
$\operatorname{SU}(2)$ are the binary
cyclic, binary dihedral, binary
tetrahedral, binary octahedral and
binary icosahedral groups.
This is the ADE list through the
McKay correspondence.

See the constant curvature discussion in
Section
\ref{%
differential-geometry-%
section-curvature%
},
especially Equation
\ref{%
differential-geometry-%
equation-constant-curvature%
}.
For the ADE side, see Section
\ref{representation-theory-section-sl2c}.

\medskip\noindent
{\bf Warning.}

The general spherical space-form
statement is $S^3/\Gamma$, with
$\Gamma<\operatorname{SO}(4)$ finite and
free.
The ADE statement is the clean
single-action case
$\Gamma<\operatorname{SU}(2)$.

\medskip\noindent
{\bf Main theorem.}

\begin{theorem}
\label{theorem-spherical-gh-boundary}
Let
$$
        S=\operatorname{SU}(2)/\Gamma
$$
be a spherical $3$-manifold, with
$\Gamma<\operatorname{SU}(2)$ finite.
Then there exists a closed orientable
hyperbolic surface $\Sigma$ and a flat
principal $\operatorname{SU}(2)$-bundle
$$
        \operatorname{SU}(2)\to P\to\Sigma
$$
such that $S$ occurs as a boundary point
of the Gromov--Hausdorff closure of
$$
        {\mathcal S}(P)
        =
        \operatorname{Met}(P)/
        \operatorname{Diff}(P).
$$
Equivalently, there are smooth metrics
$g_i$ on $P$ such that
$$
        (P,d_{g_i})\to (S,d_{g_0})
$$
in Gromov--Hausdorff distance.
\end{theorem}

Here $g_0$ is the spherical quotient
metric.
I am recording the theorem in the form
understood from the talk.
The exact statement should be checked
against the speaker's paper.

The point is that a fixed smooth
manifold $P$ can carry metrics which
collapse, in the Gromov--Hausdorff
sense, to a lower dimensional spherical
space form.

\medskip\noindent
{\bf Holonomy and monodromy.}

Let $\pi:P\to M$ be a principal
$G$-bundle with connection $\theta$.
Fix $m\in M$ and $u\in P_m$.
If $\gamma$ is a loop based at $m$,
let $\widetilde\gamma$ be its
horizontal lift starting at $u$.
Then $\widetilde\gamma(1)$ lies in
$P_m$, so there is a unique
$g_\gamma\in G$ with
$$
        \widetilde\gamma(1)
        =
        u\cdot g_\gamma .
$$

The holonomy group is
$$
        \operatorname{Hol}_u(\theta)
        =
        \{g_\gamma\}.
$$
The restricted holonomy group
$\operatorname{Hol}_u^0(\theta)$ is
obtained from null-homotopic loops.
For a connected base it is the
identity component of the holonomy
group.

The monodromy map is
$$
        m_\theta:\pi_1(M,m)\to
        \operatorname{Hol}_u(\theta)/
        \operatorname{Hol}_u^0(\theta),
        \qquad
        [\gamma]\mapsto [g_\gamma].
$$
So the quotient by restricted holonomy
records the discrete part of holonomy
seen by the fundamental group.

\medskip\noindent
{\bf Ambrose--Singer theorem.}

Let $\Omega_\theta$ be the curvature
form of $\theta$.
The Ambrose--Singer theorem says that
$$
        \mathfrak{hol}_u(\theta)
        =
        \operatorname{Lie}
        \operatorname{Hol}_u^0(\theta)
$$
is generated by curvature.
More exactly, take a point $v\in P$
reachable from $u$ by a horizontal
curve, and take horizontal vectors
$X,Y\in T_vP$.
Then the curvature value
$$
        \Omega_\theta(v)(X,Y)
        \in\mathfrak g
$$
transported back to $u$ contributes to
$\mathfrak{hol}_u(\theta)$.
The span of all such values is the
holonomy Lie algebra.

For a vector bundle connection $\nabla$,
this is the formula
$$
        \mathfrak{hol}_x(\nabla)
        =
        \operatorname{span}\{
        P_\gamma^{-1}R_y(X,Y)P_\gamma
        \}.
$$
Here $\gamma:x\to y$, $P_\gamma$ is
parallel transport, and
$X,Y\in T_yM$.

Thus curvature gives the infinitesimal
part of holonomy, while monodromy
records the residual discrete part.

\medskip\noindent
{\bf Riemann--Hilbert correspondence.}

For the smooth discussion here, the
relevant case is the flat case:
$$
        \Omega_\theta=0.
$$
Then parallel transport depends only on
the homotopy class of the path.
Choosing $m\in M$ and $u\in P_m$ gives
the holonomy representation
$$
        \rho_\theta:\pi_1(M,m)\to G,
        \qquad
        [\gamma]\mapsto g_\gamma .
$$

The Riemann--Hilbert correspondence
identifies flat principal $G$-bundles
with representations of the
fundamental group:
$$
        \left\{
        \hbox{flat principal }G
        \hbox{-bundles on }M
        \right\}/\cong
        \longleftrightarrow
        \operatorname{Hom}(\pi_1(M,m),G)/G .
$$
The group $G$ acts on representations
by conjugation.

Thus two flat bundles with connection
$(P,\theta)$ and $(Q,\omega)$ are
equivalent if there exists a principal
$G$-bundle isomorphism
$$
        \varphi:P\to Q
$$
such that
$$
        \theta=\varphi^*\omega .
$$
On the representation side this says
that the holonomy representations
$\rho$ and $\rho'$ are conjugate:
$$
        \rho'=g\rho g^{-1}.
$$
If a point in the fiber over $m$ is
fixed, the conjugacy ambiguity
disappears.

Conversely, a representation
$$
        \rho:\pi_1(M,m)\to G
$$
defines a flat principal bundle by
$$
        \widetilde M\times_\rho G .
$$
This is the flat case of the
Riemann--Hilbert dictionary.

\medskip\noindent
{\bf Collapse sequence.}

The paper uses a Kaluza--Klein type
family of metrics
$$
        g_{P,n}
        =
        {1\over n}\pi^*(h_\Sigma)
        +
        (\theta,\theta)_{
        \operatorname{SU}(2)}
$$
on the total space $P$.
As $n\to\infty$, the horizontal
directions shrink and one expects the
Gromov--Hausdorff limit to be the
spherical quotient determined by the
holonomy.

\medskip\noindent
{\bf Gromov--Hausdorff distance.}

Let $(X,d_X)$ and $(Y,d_Y)$ be compact
metric spaces.
A correspondence is a subset
$$
        R\subset X\times Y
$$
whose two coordinate projections are
surjective.

A correspondence is like a generalized
graph.
The graph of a map $f:X\to Y$ is a
correspondence only when $f$ is
surjective.
In general, $R$ behaves like a
multivalued graph.

The distortion of $R$ is
$$
        \operatorname{dis}(R)
        =
        \sup\{
        |d_X(x,x')-d_Y(y,y')|:
        (x,y),(x',y')\in R
        \}.
$$

The Gromov--Hausdorff distance is
$$
        d_{GH}(X,Y)
        =
        {1\over 2}
        \inf_R\operatorname{dis}(R),
$$
where $R$ runs through all
correspondences between $X$ and $Y$.

A small-distortion correspondence is an
almost distance-preserving matching.

\medskip\noindent
{\bf Isometry classes.}

The Gromov--Hausdorff distance is
invariant under isometry.
If
$$
        X\cong X',
        \qquad
        Y\cong Y',
$$
then
$$
        d_{GH}(X,Y)=d_{GH}(X',Y').
$$
Hence $d_{GH}$ descends to isometry
classes:
$$
        d_{GH}([X],[Y])=d_{GH}(X,Y).
$$

Moreover, for compact metric spaces,
$$
        d_{GH}([X],[Y])=0
        \quad\Longleftrightarrow\quad
        X\cong Y.
$$
So compact metric spaces modulo
isometry form a genuine metric space.

\medskip\noindent
{\bf Superspace.}

Let $P$ be a smooth manifold and
denote by $\operatorname{Met}(P)$ the
space of smooth Riemannian metrics on
$P$.
The diffeomorphism group acts by
pullback.
The superspace of $P$ is
$$
        {\mathcal S}(P)
        =
        \operatorname{Met}(P)/
        \operatorname{Diff}(P).
$$
Thus two metrics define the same point
when they differ by a diffeomorphism.

Assume now that $P$ is compact.
Each metric $g$ defines the compact
metric space $(P,d_g)$, so there is a
map
$$
        {\mathcal S}(P)\to{\mathcal{GH}},
        \qquad
        [g]\mapsto[(P,d_g)].
$$
It is well-defined because if
$g=\phi^*g'$, then
$$
        \phi:(P,d_g)\to(P,d_{g'})
$$
is an isometry.

Conversely, by Myers--Steenrod, an
isometry between smooth Riemannian
manifolds is smooth.
Thus superspace embeds in the
Gromov--Hausdorff moduli space of
compact metric spaces.

\medskip\noindent
{\bf Riemannian submersion lemma.}

A technical lemma used in the proof
gives a Riemannian submersion
$$
        f:(P,g)\to
        (\operatorname{SU}(2)/\Gamma,g_0).
$$
Here $g_0$ is the quotient of the
round metric on
$\operatorname{SU}(2)\cong S^3$.

Recall that
$f:(P,g)\to(N,h)$ is a Riemannian
submersion if $f$ is a submersion and
$$
        df_p:
        (\operatorname{Ker}df_p)^\perp
        \to T_{f(p)}N
$$
is an isometry for every $p\in P$.
Thus
$$
        T_pP
        =
        \operatorname{Ker}df_p
        \oplus
        (\operatorname{Ker}df_p)^\perp
$$
is the vertical-horizontal splitting.

For horizontal vectors $X,Y$ one has
$$
        g(X,Y)=g_0(df(X),df(Y)).
$$

This lemma is the bridge between the
principal-bundle construction and the
Gromov--Hausdorff collapse:
the ADE quotient is not only a
holonomy target, but also the metric
quotient sitting below $(P,g)$.

\medskip\noindent
{\bf Examples.}

The Poincare homology sphere is
$$
        \Sigma(2,3,5)\cong S^3/I^*
$$
where
$I^*<\operatorname{SU}(2)$ is the
binary icosahedral group of order
$120$.
It double covers the usual icosahedral
rotation group
$A_5<\operatorname{SO}(3)$.
Thus the quotient is by the binary
icosahedral group, not by the ordinary
icosahedral group.

It is a homology sphere:
$$
        H_*(\Sigma(2,3,5),\mathbb{Z})
        \cong H_*(S^3,\mathbb{Z}).
$$
but it is not simply connected.
Its fundamental group is $I^*$.

The Bolza surface is the smooth
completion of
$$
        y^2=x^5-x.
$$
It is a two-fold cover of
$\mathbb{P}^1$ branched over
$$
        0,\ \infty,\ 1,\ -1,\ i,\ -i.
$$
By Riemann--Hurwitz, a double cover of
$\mathbb{P}^1$ branched over six points has
genus
$$
        g={6-2\over 2}=2.
$$

The Bolza surface is the unique
{\it global} maximum of the systole
function on the moduli space
${\mathcal M}_2$, after giving every
Riemann surface its uniformizing
hyperbolic metric.
In other words, among closed
hyperbolic surfaces of genus $2$,
it maximizes the length of the
shortest closed geodesic.

\medskip\noindent
{\bf Points to check.}

\begin{enumerate}
\item Is the theorem stated for all
spherical 3-manifolds, or only for the
single-action quotients
$\operatorname{SU}(2)/\Gamma$?
\item In the construction, which
representation
$\rho:\pi_1(\Sigma)\to\Gamma$ is chosen?
\item What is the exact collapse
sequence $g_i$ on $P$?
\item Where is the Riemannian submersion
lemma used in the proof?
\end{enumerate}

\section{Updates on flows of
geometric structures}
\label{section-updates-on-flows-of-
geometric-structures}

\medskip\noindent
{\bf Speaker.}

Henrique Sa Earp
(Campinas, Brazil)

\medskip\noindent
{\bf Title.}

Updates on flows of geometric
structures.

\medskip\noindent
{\bf Abstract.}

Aiming at a public with interests
among Riemannian and complex
geometry,
Lie groups and parabolic PDEs,
I will recall some results towards a
general analytical theory for flows
of $H$-structures,
obtained with Fadel,
Loubeau and Moreno.
Then, as a concrete example,
I will present some recent progress
on $SU(2)$-flows on 4-manifolds,
initiating a classification
programme of ``parabolic'' flows,
based on a representation-theoretic
method originally due to Bryant
in the context of $G_2$ geometry,
obtained with Fowdar.

\medskip\noindent
{\bf Baby summary.}

The talk asks which geometric
structures can be moved by a
well-posed PDE.
First one sets up the general
$H$-structure framework.
Then one studies $SU(2)$ examples and
uses representation theory to list
the allowed flow terms.

\medskip\noindent
{\bf H-structures.}

\begin{definition}
Let $M^n$ be a smooth manifold and
let $H\subset \operatorname{GL}(n,
\mathbb{R})$.
An $H$-structure on $M$ is a
reduction of the frame bundle
$\operatorname{Fr}(M)$ from
$\operatorname{GL}(n,\mathbb{R})$
to $H$.
Equivalently, it is a section of
$\operatorname{Fr}(M)/H$.
If $\xi_0$ is a reference tensor,
then $H=\operatorname{Stab}(\xi_0)$.
\end{definition}

$$
\xymatrix{
P:=\operatorname{Fr}(M)
\ar[dd]
\ar[dr]^H
&
\\
& N:=P/H
\ar[dl]_\pi
\\
M\ar@/^-1pc/@{-->}[ur]_{\sigma}
&
}
$$
and there is an associated
tensor bundle picture
$$
        \mathcal{F}\subset
        \bigoplus \mathcal{T}^{p,q}
        \to M,
        \qquad
        \pi^*\mathcal{F}\to N.
$$
Here $TP$ splits into vertical and
horizontal parts.
Write
$$
\xymatrix{
TP \ar[dd] \ar[dr]
&
\\
&
TN:=\nu\oplus\mathcal H\ar[dl]_\pi\\
M\ar@/^-1pc/@{-->}[ur]_{\sigma}
}
$$
to record the tangent-bundle picture.

In prose: an $H$-structure is a way to
equip $M$ with extra geometric
data that is modeled on a fixed tensor
or algebraic structure with symmetry
group $H$.
The section
$\sigma:M\to P/H$ picks, at each point,
an adapted frame modulo $H$.
Equivalently, one can encode the same
information by a tensor field
$\xi$ on $M$ whose stabilizer is
exactly $H$; the bundle $P/H$ is the
space of all such local choices.

\begin{example}
\label{example-harmonic-flow}
For an $H$-structure $\sigma$ with
torsion tensor $T$, one may consider
the energy
$$
        E(\sigma)=\int_M |T|^2 .
$$
The harmonic flow is the negative
gradient flow of this energy.
Its purpose is to deform $\sigma$
toward a more special or torsion-free
structure.

Then it is natural to consider the
Ricci flow of $H$-structures.
\end{example}

\begin{definition}
\label{definition-harmonic-h-structure}
Let $H\subset \operatorname{SO}(n)$.
A harmonic $H$-structure is a
critical point of
$$
        E(\xi)
        =
        {1\over 2}
        \int_M |T|^2 .
$$
Here $\xi$ is the corresponding
$H$-structure and $T$ is its torsion
tensor.
\end{definition}

\begin{theorem}[$\varepsilon$-regularity]
\label{theorem-epsilon-regularity}
Let $(M,g)$ be closed and let
$\xi(t)$ be a harmonic $H$-flow on
$M\times[0,\tau)$.
Fix $0<r_M<\operatorname{inj}(M,g)$ and
$E_0\in(0,\infty)$.
If the local energy of $\xi(t)$ on a
geodesic ball $B_{r_M}(y)$ is small
enough, then torsion stays controlled
on a smaller ball for a short time.
\end{theorem}

\begin{theorem}[Energy gap]
\label{theorem-energy-gap}
Let $(M,g)$ be a closed Riemannian
$n$-manifold admitting a compatible
$H$-structure.
Then there is a constant
$\varepsilon(M,g,H)>0$ such that if a
compatible harmonic $H$-structure
$\xi$ satisfies
$$
        E(\xi)
        =
        \|\nabla \xi\|_{L^2(M)}^2
        <
        \varepsilon(M,g,H),
$$
then $\xi$ is torsion-free:
$$
        \nabla \xi = 0.
$$
\end{theorem}

The upshot is that the flow is being
used to drive the structure into the
small-energy regime where the
$\varepsilon$-regularity theorem
forces torsion to disappear.
So the long-time picture is meant to
end at a torsion-free limit.

One can also reverse the procedure.
Let $(M,g)$ be such that there is no
torsion-free structure in the homotopy
class.
Then the question becomes whether the
flow still has a useful limit, or
whether it must break down.

I will recall results toward a general
theory for flows of H-structures,
obtained with Fadel, Loubeau and
Moreno.
Then I will present recent results on
SU(2)-flows on 4-manifolds, starting a
classification programme for
parabolic flows.
The method is representation-theoretic
and goes back to Bryant in G2 geometry;
the recent work is joint with Fowdar.

\medskip\noindent
The key concept to construct these
flows is the Fadel--Loubeau--Moreno--
SE 24 theorem, and also the Karigiannis
08 theorem: the ``diamond action''.
In this language the evolution equation
is
$$
        \partial_t \eta = A \diamond \eta
$$
with
$$
        A\equiv A(t)=S(t)+C(t),
        \qquad
        S(t)\in S^2,
        \qquad
        C(t)\in \Lambda^2.
$$

The freedom in choosing $A$ is the
freedom to choose a differential
invariant of the $H$-structure.
Bryant's 2003 point is that, for
$H\subset \operatorname{O}(n)$, the
space $V_k(\mathfrak h)$ of $k$th order
invariants is characterized by
$$
        V_k(\mathfrak h)\oplus
        (\Lambda^1\otimes S^{k+1})
        =
        (S^2\oplus\mathfrak h^\perp)
        \otimes S^k,
$$
where
$S^k=S^k(\Lambda^1)$ and
$\mathfrak h^\perp=\Lambda^2/\mathfrak h$.
The first space
$V_1(\mathfrak h)=\Lambda^1\otimes
\mathfrak h^\perp$ is the intrinsic
torsion.
The next piece $V_2(\mathfrak h)$ is
spanned by $\nabla T$ and
$\operatorname{Rm}$,
with overlap coming from the
Bianchi-type identity
$$
        \nabla T\diamond \eta
        =
        \nabla^2\eta+\hbox{l.o.t.}
$$
so the lower-order terms are not
independent.

In the simplest case $H=\operatorname{O}(n)$,
one has
$$
        V_1(\mathfrak{so}(n))=0.
$$
Also,
$$
        V_2(\mathfrak{so}(n))
        =
        \mathbf R
        \oplus
        S^2_0(\mathbf R^n)
        \oplus W,
$$
and
$$
        \operatorname{End}(\mathbf R^n)_
        {\mathfrak{so}(n)}
        =
        \mathbf R
        \oplus
        S^2_0(\mathbf R^n).
$$

\medskip\noindent
{\bf SU(2) flows.}

Any $SU(2)$-flow can be written as
$$
        \partial_t \omega = A\diamond\omega .
$$
The point of the theory is that the
possible choices of $A$ form a finite
dimensional space of differential
invariants.
So one can classify the flows.
For $SU(2)$ this leads to a list of
quadratic torsion functionals.
There are 15 of them, built from
terms of the form
$$
        \int_M g(a_i,a_j)\,\mathrm{vol}
        \qquad\hbox{and}\qquad
        \int_M g(a_i,J_j a_k)\,
        \mathrm{vol}.
$$
Among them are the Dirichlet energy
$$
        \int_M |T|^2\,\mathrm{vol}
$$
and the complex-geometric energy
$$
        \int_M |N_J|^2\,\mathrm{vol}.
$$
So the classification problem is
really a representation-theoretic
bookkeeping problem: write down the
allowed invariant tensors, then choose
a linear combination that gives a
parabolic flow.

\medskip\noindent
In other words, the flow is not
arbitrary.
It is built from the finite list of
quadratic expressions that the
representation theory allows.
The hard part is to choose the
coefficients so that the evolution
equation is still geometric and
parabolic.
That is why the talk keeps returning
to invariant tensors, torsion, and
curvature.

\medskip\noindent
{\bf Favorite flow.}

Consider the two-parameter rich flow
$$
        \partial_t\omega=
        \Bigl(
        -2\operatorname{Ric}(g)
        -\lambda_1\mathcal{L}_{
        \sum_{i=1}^2 J_i a_i}g
        +\lambda_2\sum_{i=1}^3
        \operatorname{div}(a_i)
        \Bigr)
        \diamond\omega .
$$
It has short-time existence and
uniqueness when
$$
        -4\sqrt 5<
        \lambda_2+\lambda_1/2<0 .
$$
This is the favorite flow.

\medskip\noindent
{\bf Questions and outlook.}

Can one add terms involving $W^+$
to the flow and still keep short-time
existence?
Are there good choices of such
curvature terms?
The Lie-derivative term can be more
generally taken as
$$
        \sum_{i,j}\lambda_{ij}
        \mathcal{L}_{J_i a_j}g.
$$
What lower-order terms should one
choose?
Are there flows preserving torsion
classes, for instance Laplacian-type
flows?
What about long-time existence,
singularity formation, and solitons?
More broadly, can the community
classify the good flows for other
groups $H$?

\medskip\noindent
{\bf Audience questions.}

Is one tacitly assuming the flow is
quasi-linear?
Is that crucial?

The answer is that one wants to
control low-order derivatives of the
torsion.
Once those are under control, the
rest follows.
So the point is to control the
derivatives of the norm of the
torsion.

Would it make sense to study flows
that are not quasi-linear?
For example, one could take the
Riemann curvature tensor, contract it
in many ways, and ask what happens.

Why would one do that?
The lower-order terms depend on which
PDE one chooses.
The hope is to use this as technology
for solving geometric PDEs.
That is the perspective here.

\medskip\noindent
{\bf Baby summary.}

The talk is about moving geometric
structures by PDEs.
The structure has torsion, and the
flow tries to reduce it.
Representation theory tells you which
terms you are allowed to use.
In the best cases, the flow is
well-posed and drives the structure
toward something more rigid.
The $SU(2)$ case is especially rich,
because there are many explicit
quadratic terms to choose from.

\section{Atom applications and
MMP}
\label{section-atom-applications-and-
mmp}

\medskip\noindent
{\bf Speaker.}

Ludmil Katzarkov

\medskip\noindent
{\bf Title.}

Atom applications and MMP.

\medskip\noindent
{\bf Abstract.}

In this talk we introduce the
theory of Atoms.
Applications and connection with
MMP will be discussed.

\medskip\noindent
{\bf Birational invariants.}

\begin{enumerate}
\item LG-modules.
\item Orlov spectra.
\item Differential equations to the
monads.
\item Road map to the twistor
examples.
\end{enumerate}

\medskip\noindent
{\bf Orlov spectrum.}

Take a generator $g$ of
$D^b(X)$ and generate by taking
sums, summands, triangles and
shifts.
Let $t(g)$ be the minimum number of
triangles needed to get to any
object in $D^b(X)$.

The Orlov spectrum is the list of
generation times of
$D^b(X)$.
In the shorthand from the talk,
$$
D^b(X)=
\langle
t_1(g),\dots,t_n(g)
\rangle.
$$

\begin{enumerate}
\item $D^b(\mathbb{P}^1)$ has Orlov
spectrum $(1,2)$.
\item $D^b(E)$ has Orlov spectrum
$(1,2,2,4)$.
\end{enumerate}

Conjecture: this number is the
dimension of the variety.
This was confirmed in dimension $1$.

There are gaps in the spectrum.
A conjecture by Ballard, Favero and
Katzarkov says that if the gap is
greater than $\dim X-2$, then
$X$ is not rational.

\medskip\noindent
{\bf Differential equations.}

Now we take a shortcut and use
differential equations.
Keep in mind the Landau--Ginzburg
model, but something else will
happen.
Take a Fano manifold $X$ and
consider it as an algebraic
variety, but also as a symplectic
manifold.

\medskip\noindent
{\bf Main objective.}

If $X$ is a variety over
$\mathbb{C}$,
then birational geometry asks when
$$
\mathbb{C}(X)
\cong
\mathbb{C}(x_1,\dots,x_n),
\qquad
n=\dim X.
$$
This is the rationality problem in
its simplest form.

\medskip\noindent
{\bf Mirror symmetry.}

For smooth $X$, homological mirror
symmetry says
$$
D^b(X)
\leftrightarrow
\operatorname{FS}(Y,W),
$$
where $Y$ is an open symplectic
manifold and $W$ is a function.
The category on the right is built
from Lagrangian submanifolds, and
morphisms are given by Floer
homology.

\medskip\noindent
{\bf Example.}

For $X=\mathbb{P}^2$ one has
$$
D^b(\mathbb{P}^2)=
\langle
\mathcal{O},
\mathcal{O}(1),
\mathcal{O}(2)
\rangle.
$$
The mirror is
$$
Y=(\mathbb{C}^*)^2,
\qquad
W=x+y+\frac{1}{xy}.
$$

The observation made by them is that
the singularity theory on the
symplectic side recovers, my guess,
the birational data on the algebraic
side.

\medskip\noindent
{\bf Blow-up.}

Now turn to birational geometry.
What happens when you blow up a
point of $\mathbb{P}^2$?

$$
D^b(\operatorname{Bl}_p \mathbb{P}^2)=
\langle
\mathcal{O},
\mathcal{O}(1),
\mathcal{O}(2),
E
\rangle,
$$

and on the symplectic side we get
$I_8$
(singularity).

\medskip\noindent
{\bf Cubic threefold.}

Next example:
$X_3 \subset \mathbb{P}^4$.
The classical approach is by studying
$$
H^{2,1}(X_3)/
H_3(X_3,\mathbb{Z}),
$$
that is, the intermediate Jacobian.
This gives a moduli space called the
Theta divisor
$$
\Theta.
$$

\begin{example}
From the HMS point of view,
we have a symplectic manifold
and a potential
$$
W=
\frac{(x+y+z)^3}{xyz}+z.
$$
Some singular K3 surfaces appear.
There is also a complicated thing
at infinity.

Here one writes
$$
D^b(X_3)=
\langle
\mathcal{A},
E_1,
E_2
\rangle.
$$
\end{example}

\begin{theorem}[(-, Prizholkovsky, 2005)]
Let $X$ be an equidimensional Fano
with $\Pic X = \mathbb{Z}$,
and suppose $X$ is not
$\mathbb{P}^3$.
Consider the monodromy action on
near-by fibers of the symplectic
side.
The monodromy is a $3\times 3$
matrix.
If $\mu$ is strictly unipotent,
then $X$ is rational.
\end{theorem}

\medskip\noindent
{\bf Differential equations.}

Now we shall take a shortcut:
we use differential equations.
Keep in mind the
Landau--Ginzburg model,
but something else will happen.
Take a Fano manifold $X$ and
consider it as an algebraic
variety, but also as a symplectic
manifold.

We have de Rham cohomology, and also
a noncommutative Hodge package
$$
H_{nc}=
\bigoplus_{\lambda_i}
H_{\lambda_i}.
$$
The formula for $E_{\lambda_i}$ in
the notes is too unclear to
reconstruct exactly, so I do not
guess it here.

\begin{theorem}[Kontsevich,
Pantev, Yu]
The object
$$
\bigoplus_{\lambda_i}
H_{\lambda_i}
$$
is a birational invariant.
\end{theorem}

The quantum differential equation is
$$
\left(
\frac{\partial}{\partial u}
-\frac{K}{u^2}
+\frac{G}{u}
\right)
\psi(u)=0.
$$
where
$$
\psi=u^{\delta}g(1,\dots).
$$

\begin{theorem}
Let $X \subset \mathbb{P}^N$
be a Fano manifold of dimension
$\geq 3$.
Then
$$
\delta_{\max}
=
\dim X -1
\cdots
$$
If MMP works, and we can
...
then you might be able to conclude
that this $7$-dimensional cubic is
Picard as well.
\end{theorem}

\medskip\noindent
{\bf G-equivariant Hodge atoms.}

\medskip\noindent
{\bf Part 1.}

Review of coarse Hodge atoms.

\medskip\noindent
Let $X/\mathbb{C}$ be a smooth
projective variety of dimension at
least $2$.
For instance, one can think of a
compact K\"ahler manifold.
Let $G<\operatorname{Aut}(X)$ be a
finite subgroup.
We call $X$ a $G$-variety if $G$
acts generically freely on $X$.
The pair $(X,G)$ is non-$G$-
linearizable if there is no linear
$G$-action on $\mathbb{P}^n$ and no
$G$-equivariant birational map
between $X$ and $\mathbb{P}^n$.

\medskip\noindent
The equivariant version of the
coarse Hodge atoms keeps track of
the $G$-action on the Hodge
package.
The natural coarse version is
obtained by decomposing the atomic
package into $G$-isotypical
pieces.

\medskip\noindent
{\bf Application.}

These invariants give obstructions
to $G$-linearizability.
In particular, they produce many
examples of non-$G$-linearizable
$G$-actions on projective
varieties.

\section{Recent progresses in the
theory of atoms}
\label{section-recent-progresses-in-
the-theory-of-atoms}

\medskip\noindent
{\bf Speaker.}

Leonardo Cavenaghi

\medskip\noindent
{\bf Title.}

Recent progresses in the theory of
atoms.

\medskip\noindent
{\bf Abstract.}

In this talk we report on new
results related to the theory of
atoms obtained in collaboration
with L. Katzarkov and M.
Kontsevich. They will appear in the
text titled "9.5 lectures in the
theory of atoms", in the final
stage of preparation.

\medskip\noindent
{\bf Baby upshot.}

The point is to package the
Hodge-theoretic and quantum
information into atoms.
These atoms carry both the
representation-theoretic and
birational content.

\medskip\noindent
{\bf Relation with
Mumford--Tate groups.}

This fits the notes in
\ref{mumford-tate-%
section-mt-hodge-structures}.
There the Mumford--Tate group
controls Hodge symmetries.
The same Hodge package is also
developed in the complex-geometry
notes.
Here the Hodge atoms are a more
refined bookkeeping device:
they split the Hodge package into
pieces on which
$
\operatorname{Hod}(\mathbb{k})
$
and its equivariant variants act.

\medskip\noindent
{\bf G-equivariant atoms and
quantum cohomology.}

\medskip\noindent
{\bf Theorem.}

For any multi-degree hypersurface
$$
X:=X(1,1,1,1)\subset(\mathbb{P}^1)^4,
$$
for instance the blow-up of
$\mathbb{P}^1\times\mathbb{P}^1
\times\mathbb{P}^1\times\mathbb{P}^1$,
the invariant $\mathbb{Z}/2$-action
by swapping the first and second,
and the third and fourth factors,
is not $\mathbb{Z}/2$-linearizable.
The same holds for
$$
X\times\mathbb{P}^1,
$$
where $\mathbb{Z}/2$ acts trivially
on the $\mathbb{P}^1$ factor.

\medskip\noindent
{\bf Coarse atoms.}

For this we introduce a coarse
version of the theory of
$G$-equivariant atoms.

\medskip\noindent
{\bf Notation.}

Take your variety and its rational
cohomology
$$
V:=H^*(X,\mathbb{Q}),
$$
seen as a supervector space.
\begin{itemize}
\item Coordinates:
$$
\mathbf{t}=\{t_a\}
$$
formal dual super coordinates.
\item Novikov ring:
to track effective curve classes
$\beta \in \overline{\operatorname{NE}}
(X,\mathbb{Z})$.
The small Novikov ring is
$$
\mathcal{N}_X.
$$
\noindent
It is the ring of power series.
\item Coefficient ring:
formal power series.
\end{itemize}

The Gromov--Witten potential allows
us to define a new product in this
ring that recovers the classical
Poincar\'e cup product.
\begin{itemize}
\item Classical part:
intersection theory.
\item Quantum correction:
curve counting, enumerative
geometry.
\end{itemize}

\medskip\noindent
{\bf Big quantum cohomology.}

\begin{itemize}
\item Space:
$$
V_{\mathrm{big}}=
V\otimes_{\mathbb{Q}}\mathcal{N}_X.
$$
\item Product: third-order
differential operator.
\item Deformation: full deformation
in $\mathbf{t}$.
\end{itemize}

\medskip\noindent
{\bf Small quantum cohomology.}

\begin{itemize}
\item Space:
$$
V\otimes\mathcal{N}_X.
$$
\item Product:
[missing]
\item Deformation:
[missing]
\end{itemize}

\medskip\noindent
{\bf Example.}

For $X=\mathbb{P}^2$ with basis
$$
\{\Delta_0,\Delta_1,\Delta_2\}
=\{1,H,P\},
$$
the Gromov--Witten potential is
$$
\phi_X(\mathbf{t},q)=\cdots
$$
and it has a positive radius of
convergence.

Big quantum product relations.

\medskip\noindent
{\bf Questions.}

\begin{enumerate}
\item Analytic germ.
\item Tangent space product.
\item Logarithmic tangent bundle.
\end{enumerate}

\medskip\noindent
{\bf General logarithmic tangent
bundle.}

Idea: to have a germ of a manifold
such that the tangent space is the
cohomology vector space and the
quantum product is an associative
product in this vector space.

For $X/\mathbb{C}$ smooth
projective, let
$$
V:=H^*(X,\mathbb{Q}).
$$
Let $\Phi_X$ be the
Gromov--Witten potential.
Then the third derivatives
$\partial^3\Phi_X$ yield a unital,
commutative, associative product on
the tangent bundle of a germ
$$
M\subset H^*(X,\mathbb{C}).
$$
This is the Frobenius-manifold
picture behind quantum cohomology.

\medskip\noindent
The problem is that convergence is
most of the time hypothetical.
The strategy is to introduce an
auxiliary non-Archimedean field
$\mathbb{k}$ of characteristic $0$
to ensure convergence.
The idea is to replace the germ
$M$ by a manifold $B_X$ where we do
have convergence for the
Gromov--Witten invariants.
More precisely, for $b \in B_X$,
$$
T_bB_X\simeq H^*(X,\mathbb{k},*_b).
$$

\begin{definition}
The manifold $B_X$ is called an
$A$-model manifold.
\end{definition}

Now we put an Euler vector field
on $B_X$.
If $*$ is a tensor field on $B_X$
which is actually a Higgs field,
then $\operatorname{Eu}$ scales the
product.
That is,
$$
\mathcal{L}_{\operatorname{Eu}}(*)=*.
$$
Equivalently, we have an
automorphism
$$
k_b:=\operatorname{End}_{\operatorname{Eu}}
(*_b)
\in \operatorname{End}
\bigl(H^*(X,\mathbb{k})\bigr).
$$
That is the starting point of atom
theory, from our perspective.
There are other starting points.

\medskip\noindent
{\bf Disjoint unions.}

Thom--Sebastiani principle.
If a variety is a disjoint union
$$
X\sqcup Y,
$$
its Gromov--Witten potential
uncouples as a sum of the factors:
$$
\Phi_{X\sqcup Y}
=
\Phi_X \oplus \Phi_Y.
$$
This is important because we need
to study blow-ups, the main
operations in birational geometry.

\begin{theorem}[Iritani, 2003]
Let $\widetilde X$ be the blow-up
of $X$ along a smooth subvariety
$Z\subset X$ of codimension
$r\geq 2$.
There exists a canonical analytic
identification between open
domains
$$
\psi:
B_{\widetilde X}\supset U
\xrightarrow{\sim}
V\subset
B_X\times B_Z\times\cdots\times B_Z.
$$
Moreover, one can choose the open
domains so that the Euler vector
field and the product are
preserved:
$$
\psi^*(*)=*,
\qquad
\psi^*(\operatorname{Eu})
=\operatorname{Eu}.
$$
\end{theorem}

\medskip\noindent
{\bf Hodge action and fixed locus.}

$\operatorname{Hod}(\mathbb{k})$
acts on $B_X$ preserving the quantum
product $*$ and the Euler field
$\operatorname{Eu}$.

On coordinates:
\begin{itemize}
\item $U_1$ (N\"eron--Severi):
trivial action, corresponding to
algebraic classes.
\item $U_2, U_{\mathrm{odd}}$:
scaling by Hodge weights.
\end{itemize}

\begin{definition}
The Hodge fixed locus
$$
B_X^{\operatorname{Hod}}\subset B_X
$$
is the subvariety where
$\operatorname{Hod}(\mathbb{k})$
acts trivially.
Geometrically, it is the locus
where all non-$(p,p)$ coordinates
vanish.
\end{definition}

$$
\mathcal{HS}\simeq\operatorname{Rep}
\operatorname{Hod}
$$
for $\mathbb{Q}$-polarizable pure
Hodge structures.

We can decompose the tangent space
to $B_X$ using generalized
eigenspaces:
$$
T_bB_X=
\bigoplus_{\lambda\in\operatorname{Spec}
(k_b)}V_\lambda.
$$
The catch is that the group is
pro-reductive.
This means we can write, in fact,
$$
T_bB_X=
\bigoplus_{\lambda\in\operatorname{Spec}
(k_b)}
\bigoplus_{\gamma}
V_\lambda\cap V_\gamma.
$$
The summands
$$
\bigoplus_{\gamma}
V_\lambda\cap V_\gamma
$$
are called {\it coarse Hodge atoms}.

Each atom $\alpha$ corresponds to a
$\overline{\mathbb{Q}}$-linear
representation
$$
E^\alpha
$$
of
$\operatorname{Hod}_{\overline{\mathbb{Q}}}$,
independent of the point $b$.

\medskip\noindent
{\bf Numerical invariants.}

The Hochschild grading $(p-q)$ on
$E^\alpha\otimes_{\overline{\mathbb{Q}}}
\mathbb{C}$ defines:
\begin{enumerate}
\item Rank of Hodge classes:
$$
\rho_\alpha
=
\dim
\left(V_\lambda\cap H_{p,q}\right).
$$
\item Hodge polynomial.
It is characterized by the
coefficient:
$$
P_\lambda(t)=
\dim\left(
V_\lambda\cap H^{p-q=n}(X)
\right).
$$
\end{enumerate}

\begin{theorem}[KKPY
Non-rationality criterion]
Let $X$ be an $n$-dimensional
smooth complex projective variety.
If there is a coarse Hodge atom in
the Hodge atomic composition of $X$
that does not appear in the Hodge
atomic composition of any variety
of dimension $\leq n-2$, then $X$
is not rational.
\end{theorem}

\medskip\noindent
{\bf Presence of $G$-actions.}

The action on the variety defines
an action on the cohomology.
Now one does not look at
$\operatorname{Hod}(\mathbb{k})$
alone, but at
$$
\operatorname{Hod}(\mathbb{k})
\times G,
$$
which acts on $B_X$.
Restricting to
$$
B_X^{G\times\operatorname{Hod}},
$$
a coarse $G$-atom $\alpha$ is the
isomorphism class of a
$G\times\operatorname{Hod}(\mathbb{k})$-
representation on
$\mathcal{B}_{\lambda,b}$.

\medskip\noindent
{\bf $G$-invariant Hodge classes.}

$$
\rho_\alpha^G:=
\dim_{\mathbb{C}}
\left(
E^\alpha\otimes_{\overline{\mathbb{Q}}}
\mathbb{C}
\right)^{
G\times\operatorname{Hod}
}.
$$

\medskip\noindent
{\bf Equivariant constraint.}

Since the algebraic unit $1$ is
always $G\times\operatorname{Hod}$
invariant, we have
$$
\rho_\alpha^G\geq 1.
$$

\medskip\noindent
{\bf Obstructions to
$G$-linearizable actions.}

\medskip\noindent
{\bf Atomic constraints in low
dimensions ($\leq 1$).}

Atoms originating from points
$(G/H)$ or curves
$(G\times_H C)$ satisfy:
\begin{itemize}
\item Rank bound:
$$
\rho_\alpha^G\in\{1,2\}.
$$
\item Rigidity:
if $\rho_\alpha^G$ is constant,
there is only one invariant Hodge
class.
\end{itemize}

\medskip\noindent
{\bf Example.}

The hypersurface
$$
X(1,1,1,1)\subset(\mathbb{P}^1)^4.
$$
Rationality:
$X$ is rational, the blow-up of
$\mathbb{P}^1\times\mathbb{P}^1
\times\mathbb{P}^1$
along an elliptic curve $E$.
The $G$-action is
$G=\mathbb{Z}/2$, swapping factors
1 $\leftrightarrow$ 2 and
3 $\leftrightarrow$ 4.

The atomic decomposition of $X$
yields three primary coarse
$G$-atoms:
$$
\begin{array}{c|c|c|c}
\lambda & \rho & \rho^G & P_\alpha(t)\\
16 & 1 & 1 & 1\\
4 & 4 & 2 & 4\\
0 & 5 & 3 & t+3+t^{-1}
\end{array}
$$
Note that
$$
\rho_{\mathrm{total}}^G=1+2+3=6,
$$
tracking the $G$-invariant Hodge
classes.

The atom of the elliptic curve
must contribute.
It contributes with:
\begin{itemize}
\item points $(1,1)$ or $(2,1)$.
\item curve $E$:
$(2,2)$.
\end{itemize}

Target atom $\lambda=0$:
$$
(\rho,\rho^G)=(5,3).
$$

If the action was linearizable,
its atoms must be constructible
from atoms of $\mathbb{P}^3$ and
$G$-invariant subvarieties of
dimension $\leq 1$.

The contradiction:
...

\begin{theorem}[Rationality criterion]
If $X$ is a smooth complex
projective variety of dimension
$n\geq 2$, and the maximal generic
spectral dimension on the Hodge
locus $B_X^{\operatorname{Hod}}$
satisfies
$$
\mu(X)>n-2,
$$
then $X$ is not rational.
\end{theorem}

\medskip\noindent
{\bf Quantum differential equation.}

Over the non-Archimedean field,
recall that we have a manifold and
then we put the quantum
differential equation.
We solve it via an Ansatz.
After a big analysis, we get for
the cubic, indeed,
$$
\frac{5}{3}.
$$

\end{document}
